% ============================================================
\chapter{Maximum Principle}
\label{ch:maximum}
% ============================================================

\section{Introduction}

Imagine a metal plate whose edges are heated. Where is the hottest point inside? Physical intuition is unequivocal: in the absence of an internal heat source, the temperature maximum must be attained on the boundary. This observation, elevated to the status of a mathematical principle by Eberhard Hopf in the 1920s--1930s, has become one of the most powerful tools in PDE theory. The maximum principle requires no explicit knowledge of the solution: it provides a priori bounds, uniqueness results, and comparison theorems through the sheer force of the equation itself. Hopf's lemma, which refines the principle by showing that the normal derivative is strictly negative at the maximum point, is a gem of analysis whose applications irrigate the entire elliptic and parabolic theory.

\begin{intuition}
Physically, the maximum principle expresses the fact that in the absence
of an internal heat source, the maximum temperature in a body is attained
on the boundary. More generally, an equilibrium state cannot create a new
extremum in the interior.
\end{intuition}

\section{Maximum principle for elliptic operators}

\subsection{General elliptic operators}

\begin{definition}[Elliptic operator]
A second-order differential operator:
\begin{equation}\label{eq:elliptic_operator}
    Lu = -\sum_{i,j=1}^{n} a_{ij}(x)\, u_{x_i x_j}
    + \sum_{i=1}^{n} b_i(x)\, u_{x_i} + c(x)\, u
\end{equation}
is called \textbf{elliptic} in $\Omega$ if the matrix $(a_{ij}(x))$ is
positive definite for all $x \in \Omega$, i.e., there exists $\theta > 0$
such that:
\[
    \sum_{i,j} a_{ij}(x)\, \xi_i\, \xi_j \geq \theta\, \abs{\xi}^2
    \quad \forall x \in \Omega, \; \forall \xi \in \R^n.
\]
\end{definition}

\subsection{Weak maximum principle}

\begin{theorem}[Weak maximum principle --- case $c = 0$]
\label{thm:weak_max_elliptic}
Let $\Omega \subset \R^n$ be a bounded open set and $L$ an elliptic
operator with $c = 0$. If $u \in C^2(\Omega) \cap C(\overline{\Omega})$
satisfies $Lu \leq 0$ in $\Omega$, then:
\begin{equation}\label{eq:weak_max}
    \max_{\overline{\Omega}} u = \max_{\partial\Omega} u.
\end{equation}
\end{theorem}

\begin{proof}
\textbf{Case $Lu < 0$.} Suppose $u$ attains its maximum at
$x_0 \in \Omega$. Then $\nabla u(x_0) = 0$ and $D^2u(x_0) \leq 0$
(negative semi-definite Hessian). Therefore:
\[
    Lu(x_0) = -\sum_{i,j} a_{ij}(x_0)\, u_{x_ix_j}(x_0) \geq 0,
\]
since $\sum a_{ij} \xi_i\xi_j \geq 0$ for all $\xi$ implies
$\mathrm{tr}(A \cdot D^2u) \leq 0$, hence $-\mathrm{tr}(A \cdot D^2u) \geq 0$.
This contradicts $Lu < 0$.

\textbf{Case $Lu \leq 0$.} Consider
$v_\varepsilon(x) = u(x) + \varepsilon e^{\alpha x_1}$
for $\alpha$ large enough (so that $L(e^{\alpha x_1}) < 0$). Then
$Lv_\varepsilon < 0$ and the previous case applies. Conclude by letting
$\varepsilon \to 0$.
\end{proof}

\begin{theorem}[Weak maximum principle --- case $c \geq 0$]
If $c \geq 0$ and $Lu \leq 0$ in $\Omega$, then:
\begin{equation}\label{eq:weak_max_c}
    \max_{\overline{\Omega}} u \leq \max_{\partial\Omega} u^+,
\end{equation}
where $u^+ = \max(u, 0)$.
\end{theorem}

\subsection{Strong maximum principle}

\begin{theorem}[Strong maximum principle --- Hopf]
\label{thm:strong_max}
Let $\Omega$ be a bounded connected open set and $L$ elliptic with $c = 0$.
If $u \in C^2(\Omega) \cap C(\overline{\Omega})$ satisfies $Lu \leq 0$
in $\Omega$ and attains its maximum at an interior point $x_0 \in \Omega$,
then $u$ is constant in $\Omega$.
\end{theorem}

\begin{theorem}[Hopf's lemma]
\label{thm:hopf_lemma}
Under the hypotheses of the previous theorem, if the maximum $M$ is
attained at a point $x_0 \in \partial\Omega$ (and not in the interior),
and if $\Omega$ satisfies the \textbf{interior sphere condition} at $x_0$,
then:
\begin{equation}\label{eq:hopf_lemma}
    \frac{\partial u}{\partial \nu}(x_0) > 0,
\end{equation}
where $\nu$ is the outward normal.
\end{theorem}

\begin{proof}[Proof sketch of Hopf's lemma]
One constructs a \textbf{barrier} $w(x) = e^{-\alpha\abs{x-y}^2} - e^{-\alpha R^2}$
in the annulus $B(y,R) \setminus B(y,\rho)$, where $B(y,R)$ is the interior
sphere tangent at $x_0$. For $\alpha$ large enough, $Lw < 0$. The comparison
principle gives $u - M \leq -\varepsilon w$, whence
$\frac{\partial u}{\partial \nu}(x_0) \geq \varepsilon \frac{\partial w}{\partial \nu}(x_0) > 0$.
\end{proof}

\section{Applications of the maximum principle}

\subsection{Uniqueness}

\begin{corollary}[Uniqueness for the Dirichlet problem]
The problem $Lu = f$ in $\Omega$, $u = g$ on $\partial\Omega$
(with $c \geq 0$ and $L$ elliptic) admits at most one classical solution.
\end{corollary}

\subsection{A priori estimates}

\begin{theorem}[A priori estimate]
Let $u \in C^2(\Omega) \cap C(\overline{\Omega})$ solve $Lu = f$ in
$\Omega$, $u = g$ on $\partial\Omega$, with $L$ elliptic and $c = 0$.
Then:
\begin{equation}\label{eq:a_priori_estimate}
    \max_{\overline{\Omega}} \abs{u} \leq \max_{\partial\Omega} \abs{g}
    + C\, \max_{\overline{\Omega}} \abs{f},
\end{equation}
where $C$ depends on $\Omega$, the coefficients of $L$, and $\theta$.
\end{theorem}

\subsection{Comparison principle}

\begin{theorem}[Comparison principle]
If $Lu \leq Lv$ in $\Omega$ and $u \leq v$ on $\partial\Omega$
(with $c \geq 0$), then $u \leq v$ in $\overline{\Omega}$.
\end{theorem}

\begin{proof}
Set $w = u - v$. Then $Lw \leq 0$ and $w \leq 0$ on $\partial\Omega$.
The maximum principle gives $w \leq 0$ in $\overline{\Omega}$.
\end{proof}

\section{Maximum principle for parabolic equations}

\begin{definition}[Parabolic operator]
The operator $\mathcal{L}u = u_t - Lu$ where $L$ is elliptic is called
\textbf{parabolic}.
\end{definition}

\begin{theorem}[Weak parabolic maximum principle]
Let $\Omega_T = \Omega \times (0,T]$ and $\Gamma_T$ the parabolic boundary.
If $u \in C^{2,1}(\Omega_T) \cap C(\overline{\Omega_T})$ satisfies
$\mathcal{L}u \leq 0$ in $\Omega_T$ (with $c = 0$), then:
\begin{equation}\label{eq:parabolic_max}
    \max_{\overline{\Omega_T}} u = \max_{\Gamma_T} u.
\end{equation}
\end{theorem}

\begin{theorem}[Strong parabolic maximum principle]
Under the above hypotheses, if $u$ attains its maximum at a point
$(x_0, t_0) \in \Omega_T$ with $t_0 > 0$, then $u$ is constant in
$\Omega \times [0, t_0]$.
\end{theorem}

\begin{remark}
The fundamental difference from the elliptic case is that the parabolic
maximum can only be attained on the parabolic boundary $\Gamma_T$
(which does not include $t = T$). Time plays an asymmetric role:
information propagates only into the future.
\end{remark}

\section{Sub- and supersolution method}

\begin{definition}[Subsolution, supersolution]
A function $\underline{u} \in C^2(\Omega) \cap C(\overline{\Omega})$ is a
\textbf{subsolution} of $Lu = f$ if $L\underline{u} \leq f$ in $\Omega$
and $\underline{u} \leq g$ on $\partial\Omega$. A \textbf{supersolution}
$\overline{u}$ satisfies the reversed inequalities.
\end{definition}

\begin{theorem}[Enclosure by sub/supersolutions]
If $\underline{u}$ is a subsolution and $\overline{u}$ a supersolution,
then $\underline{u} \leq u \leq \overline{u}$ in $\overline{\Omega}$.
If existence is established, the solution $u$ is trapped between
$\underline{u}$ and $\overline{u}$.
\end{theorem}

\section{Maximum principle for nonlinear equations}

\begin{theorem}[Comparison principle for $-\Delta u = f(u)$]
Let $f : \R \to \R$ be locally Lipschitz. If $-\Delta u \leq f(u)$ and
$-\Delta v \geq f(v)$ in $\Omega$, and $u \leq v$ on $\partial\Omega$,
then $u \leq v$ in $\Omega$, provided $f$ is decreasing (or one has an
a priori $L^\infty$ estimate).
\end{theorem}

\section{Exercises}

\begin{exercise}
Let $u$ be harmonic in $B(0,1) \subset \R^2$ with $u = x^2$ on
$\partial B(0,1)$. Show that $\abs{u(x)} \leq 1$ in $B(0,1)$.
Determine $u$.
\end{exercise}

\begin{exercise}
Let $Lu = -\Delta u + u$ ($c = 1 > 0$). Show that if $Lu \leq 0$ and
$u \leq 0$ on $\partial\Omega$, then $u \leq 0$ in $\Omega$.
\end{exercise}

\begin{exercise}
Use the maximum principle to show that the solution of
$-\Delta u = 1$ in $B(0,R)$ with $u = 0$ on $\partial B$ satisfies
$0 \leq u \leq R^2/(2n)$.
\end{exercise}

\begin{exercise}[Hopf's lemma]
Let $u$ be harmonic in $B(0,1)$ with $u < 1$ in $B(0,1)$ and
$u(e_1) = 1$ where $e_1 = (1,0,\ldots,0)$. Show that
$\frac{\partial u}{\partial r}(e_1) > 0$.
\end{exercise}

\begin{exercise}
Let $u$ solve $u_t = u_{xx} + u$ on $(0,\pi) \times (0,T)$ with
$u(0,t) = u(\pi,t) = 0$. Does the maximum principle apply directly?
How can it be adapted?
\emph{Hint:} set $v = e^{-t} u$.
\end{exercise}

\begin{exercise}
Show by the comparison principle that the solution of
$u_t = u_{xx}$, $u(0,t)=u(1,t)=0$, $u(x,0) = \sin(\pi x)$
satisfies $0 \leq u(x,t) \leq e^{-\pi^2 t}$ for all $t > 0$.
\end{exercise}

\begin{exercise}
Prove Liouville's theorem for harmonic functions using the maximum
principle (without the gradient estimate).
\emph{Hint:} consider $v_R(x) = u(x) - u(0) - \varepsilon(R^2 - \abs{x}^2)$.
\end{exercise}

\begin{exercise}
Let $\Omega$ be a bounded open subset of $\R^n$ and
$u \in C^2(\Omega) \cap C(\overline{\Omega})$ satisfying
$-\Delta u + c(x) u = f$ with $c \geq 0$, $f \leq 0$, and $u = 0$
on $\partial\Omega$. Show that $u \leq 0$ in $\Omega$.
\end{exercise}
